% ============================================================================
% Figure legend for Fig_LandscapeContinuum (svg / pdf / png).
% Final manuscript numbering: Fig. 3, placed in Section IV. Pure
% conceptual schematic: no simulation output, no data dependency; the
% generator is code/make_fig_landscape.py. On-chart text is
% identity-only; all interpretation lives in this legend. Format follows
% the manuscript convention ("Fig. 3 | Title." with lowercase bold panel
% letters a, b).
% ============================================================================

\noindent\textbf{Fig.~3\,|\,The continuum of module integrity: landscape
simplification and topological abstraction.}

\smallskip
\noindent\textbf{a}, Topological simplification of the attractor
landscape as module integrity declines. The four panels (from left to
right) depict a single continuum: at high integrity, the landscape
features multiple broad, shallow basins connected by open paths; with
declining integrity, the number of accessible basins decreases, their
depths increase, and the barriers between them rise; at very low
integrity, the landscape simplifies to a single deep basin, with paths
to other states being topologically blocked. Orange dots mark the
energy minima of the basins, and contour density illustrates the local
gradients.

\smallskip
\noindent\textbf{b}, Abstract topological representation of the
landscapes in panel a, with nodes spatially mapped to the corresponding
basins. From left to right, the number of accessible states decreases
and connections progressively disappear. To reflect the geometry of the
basins, each node is drawn as a solid coloured core with a dashed ring
whose radius follows the basin width in panel a, and an orange central
dot marks the energy minimum. In the high-integrity state, curved edges
illustrate the relative freedom of drift between multiple equivalent
states; as integrity declines, pathways straighten and vanish; at the
low-integrity state, a red crossed dotted line explicitly denotes the
topological blocking of the return path; at very low integrity, only
the single deep basin remains. Panels a and b share the same integrity
gradient and spatial alignment.